\documentclass{article}
\usepackage[utf8]{inputenc}
\usepackage{amsmath}
\usepackage{amssymb}
\usepackage{hyperref}
\usepackage{geometry}
\geometry{
 a4paper,
 margin=1in,
}
\begin{document}


\section{What to do with seasonality}

\subsection{What they did in the paper}
Zimmer \& Askitas (2009) explicitly used seasonally unadjusted unemployment rates and then modeled seasonality via 12-month differences instead of month dummies. They wrote:
“To calculate the change of the variables used, a 12-month lag operator is used; consequently, the past stock variables are of lag 12. This has the advantage that we do not need to model seasonality explicitly.”

\subsection{If I do it }
Use unadjusted unemployment rates (as in the paper) — and handle seasonality either by:
\begin{itemize}
    \item including month dummies in lagmodels (your current plan), or
    \item using YoY differences as a robustness check
\end{itemize}

\subsection{Why is it better}
Seasonally adjusted series often have artificial kinks or irregular first differences because the adjustment filters (X-13ARIMA-SEATS) can introduce short-term distortions.
With unadjusted data + month dummies, you keep the original data-generating process, which is crucial for interpretability and comparison with the literature


\section{Cointegration}

\subsection{What about the sample size}
Your current sample (2016–2021, = 60 months) might be too short for reliable cointegration testing (Engle–Granger) — those tests are notoriously underpowered with <100 obs.
You found no cointegration, meaning you can’t justify an ECM and should switch to short-run forecasting models in differences — which you already plan to do.


\section{Sample size}

\subsection{Ideas}:
Use at least 8–10 years (= 2011,5–2023) for your estimation sample.
This keeps enough variation for ADF to detect I(1) properly and gives 120 obs for robust forecasting.

\textbf{What does this timeframe entail}
\begin{itemize}
    \item 11 years ties 12 months = 132 observations → enough for ADF, Granger, and rolling forecasts.
    \item Covers several macro regimes:
    \begin{itemize}
        \item pre-Euro-crisis recovery
        \item COVID shock (2020–21)
        \item post-COVID labor tightening (2022–23)
    \end{itemize}
    \item  The structural breaks (COVID especially) are economically interpretable — not random noise.
\end{itemize}

\textbf{What would be caveats}
\begin{itemize}
    \item There is a structural break in 2020 (you see it in your graph).
That means: parameters before and after COVID might differ.
\item f you estimated one static model over the whole period, its coefficients could be biased or unstable.
\end{itemize}

\subsection{The first idea}
Dataset should start 2011-05-31 (7.0) and end 2023-01-31 (5.7)
That’s 11 years + 9 months = 141 months. This means all your final monthly series (unemployment + aggregated Google indicators) must exactly cover: 2011-05 → 2023-01

\section{Forecasting}
\subsection{HOw to handle the data}
Use the full 2011–2023 data for estimation + evaluation.
Model in differences to ensure stationarity.
You can treat covid as a regime change via:
\begin{itemize}
    \item Dummy variable (optional robustness)
    \item Or better: rolling-window forecasts, e.g. 60-month window, expanding/rolling origin.
\end{itemize}



\section{Google trends}

\subsection{Problem}
Google Trends allows a maximum of 5 years per request for weekly data.
Beyond that, it automatically switches to monthly data, and each query is normalized independently (each 0–100 within that time slice).
To maintain comparability, you’ll need to download overlapping 5-year windows and later rescale them.

\subsection{Stitching and rescaling Google Trends.}
Google Trends returns weekly indices scaled independently within each query window.
To obtain a consistent weekly series for 2011--05 to 2023--01, we download three overlapping
windows of up to five years each (to preserve weekly resolution). Let $W_1$, $W_2$, and $W_3$
denote the weekly series for a given keyword from the first, second, and third window, respectively.
Because each window is normalized to $[0,100]$ within the window, we rescale across overlaps.

\begin{enumerate}
\item \textbf{Rescale $W_2$ to $W_1$.} Over the common overlap $\mathcal{T}_{12}$, compute the robust
scaling factor as the median ratio
\[
s_{2\to 1} \;=\; \operatorname{median}\!\left( \frac{W_1(t)}{W_2(t)} \right)_{t \in \mathcal{T}_{12}},
\]
after excluding zeros and missing values. Multiply all observations in $W_2$ by $s_{2\to 1}$.

\item \textbf{Rescale $W_3$ to (rescaled) $W_2$.} Over the overlap $\mathcal{T}_{23}$, compute
\[
s_{3\to 2} \;=\; \operatorname{median}\!\left( \frac{W_2^{\,\ast}(t)}{W_3(t)} \right)_{t \in \mathcal{T}_{23}},
\]
where $W_2^{\,\ast} = s_{2\to 1}\,W_2$ is the rescaled $W_2$. Multiply all observations in $W_3$ by $s_{3\to 2}$.

\item \textbf{Concatenate.} Form the stitched series by taking $W_1$, then appending the rescaled $W_2^{\,\ast}$
excluding its overlap with $W_1$, and finally the rescaled $W_3^{\,\ast}$ excluding its overlap with
$W_1 \cup W_2^{\,\ast}$.

\item \textbf{Trim and aggregate.} Trim the resulting weekly series to the target window
\mbox{2011-05-01 to 2023-01-31}. For monthly models, aggregate by calendar-month averaging.
For the W34/W12 nowcasting design, compute bi-weekly means aligned to $W34(M-1)$ and $W12(M)$
and map them to month $M$.
\end{enumerate}

\subsection{Notes after stiching}
Seeing values >100 after stitching is normal. You’ve rescaled each window onto a common scale; the combined index is no longer bounded by 100.
Your printed tail looks great. The weekly window ending 2023-01-29 will aggregate into a January 2023 monthly mean (date labeled 2023-01-31) — exactly what you want for merging with the unemployment rate.

\subsection{Biweekly aggregation}

Following Zimmer et al.\ (2020), we separate each calendar month $M$ into two
half-month periods. Weeks with days $1\text{–}15$ of month $M$ form the
\emph{early-month} average $W12(M)$, while weeks with days $16\text{–}\text{end}$
of month $M{-}1$ form the \emph{late-month} average $W34(M{-}1)$, which is then
shifted forward to align with month $M$. The resulting feature panel provides
two biweekly signals per month aligned to the same target month index.


\end{document}